% ============================================================
% 4. EXPERIMENTAL SETUP
% ============================================================
\section{Experimental Setup}
\label{sec:experiments}

\subsection{Dataset}
\label{sec:dataset}

We use the Brain Tumor Detection dataset~\cite{bhuvaji2020brain}, a publicly available collection of 2D brain MRI slices from Kaggle. The dataset comprises approximately 1,500 healthy and 1,500 anomalous (tumor-bearing) JPEG images of varying original dimensions. Ground truth tumor annotations are provided as VIA polygon masks in JSON format for a subset of anomalous images.

All images are preprocessed to $256 \times 256$ grayscale PNG format and normalized to the range $[-1, 1]$. The dataset is split into training (70\%), validation (15\%), and test (15\%) partitions. The test set contains 225 healthy and 450 anomalous images. Importantly, the DDPM is trained \textit{exclusively} on healthy images (via their precomputed latent representations), following the unsupervised anomaly detection paradigm.

\subsection{Implementation Details}
\label{sec:implementation}

The pipeline is implemented in PyTorch 2.12 with the MONAI framework~\cite{monai2022} for the VAE component. All experiments are conducted on a single NVIDIA RTX 4060 GPU with 8~GB VRAM. To accommodate the memory constraints, we employ several optimizations:

\begin{itemize}[leftmargin=*]
  \item \textbf{Gradient checkpointing} in the U-Net to trade computation for memory during training.
  \item \textbf{Automatic Mixed Precision (AMP)} for both VAE and DDPM training.
  \item \textbf{Zero DataLoader workers} (\texttt{num\_workers=0}) to avoid memory duplication, as precomputed latents are already loaded into RAM.
  \item \textbf{Consolidated latent storage}: all training latents are stored in a single tensor file ($\sim$17~MB) rather than individual files, eliminating file handle overhead.
\end{itemize}

Table~\ref{tab:hyperparameters} summarizes the key hyperparameters for both training stages.

\begin{table}[H]
\centering
\caption{Training hyperparameters for the two-stage pipeline.}
\label{tab:hyperparameters}
\begin{tabular}{@{}lcc@{}}
\toprule
\textbf{Parameter} & \textbf{VAE (Stage 1)} & \textbf{DDPM (Stage 2)} \\
\midrule
Input resolution    & $256 \times 256 \times 1$ & $32 \times 32 \times 4$ \\
Channels            & $[64, 128, 256, 256]$ & $[128, 256, 512, 512]$ \\
Residual blocks     & 2 per level & 2 per level \\
Attention           & Levels 3--4 & Res.\ $16 \times 16$, $8 \times 8$ \\
Optimizer           & AdamW & AdamW \\
Learning rate       & $1 \times 10^{-4}$ & $2 \times 10^{-4}$ \\
Weight decay        & 0.01 & 0.01 \\
Batch size          & 4 (eff.\ 32) & 32 \\
Training duration   & 150 epochs & 200K steps \\
EMA decay           & 0.999 & 0.9999 \\
Gradient clipping   & 1.0 & 1.0 \\
Mixed precision     & \checkmark & \checkmark \\
\bottomrule
\end{tabular}
\end{table}

\subsection{Evaluation Metrics}
\label{sec:metrics}

\paragraph{Anomaly Detection.} We evaluate anomaly detection at two levels:
\begin{itemize}[leftmargin=*]
  \item \textbf{Image-level}: Each image receives a scalar anomaly score (mean raw L1 reconstruction error). We report the Area Under the Receiver Operating Characteristic curve (AUROC), Area Under the Precision--Recall Curve (AUPRC), F1-score at the optimal threshold, sensitivity (true positive rate), and specificity (true negative rate).
  \item \textbf{Pixel-level}: Using ground truth polygon masks, we compute pixel-level AUROC by treating each pixel's anomaly map value as a prediction and the binary mask as ground truth.
\end{itemize}

\paragraph{Synthetic Image Quality.} We assess generation quality using four complementary metrics:
\begin{itemize}[leftmargin=*]
  \item \textbf{Fr\'{e}chet Inception Distance (FID)}~\cite{heusel2017gans}: measures the distributional distance between real and generated image feature distributions.
  \item \textbf{Structural Similarity Index (SSIM)}~\cite{wang2004image}: quantifies structural similarity between randomly paired real and synthetic images.
  \item \textbf{Learned Perceptual Image Patch Similarity (LPIPS)}~\cite{zhang2018perceptual}: measures perceptual distance using deep features.
  \item \textbf{Intra-set Diversity}: LPIPS computed between pairs of generated images, measuring the diversity of synthetic outputs.
\end{itemize}

Note that SSIM and LPIPS are computed as distribution-level statistics over 200 randomly sampled cross-distribution pairs, rather than paired comparisons, since synthetic images are not reconstructions of specific real images.
